\par\indent\SourceLeaf{246}%
\textsc{Theorema.} \textit{Ductis in superficie curva ab eodem puncto initiali innumeris lineis brevissimis aequalis longitudinis, linea earum extremitates iungens ad illas singulas erit normalis.}

Operae pretium esse duximus, hoc theorema e proprietate fundamentali linearum brevissimarum deducere: ceterum eius veritas etiam absque calculo per sequens ratiocinium intelligi potest. Sint \(AB,AB'\) duae lineae brevissimae eiusdem longitudinis, angulum infinite parvum ad \(A\) includentes, supponamusque, alterutrum angulorum elementi \(BB'\) cum lineis \(BA,B'A\) differre quantitate finita ab angulo recto, unde per legem continuitatis alter maior alter minor erit angulo recto. Supponamus, angulum ad \(B\) esse \(=90^0-\omega\), capiamusque in linea \(BA\) punctum \(C\), ita ut sit \(BC=BB'\mathord{.}\mathrm{cosec}\,\omega\): hinc quum triangulum infinite parvum \(BB'C\) tamquam planum tractare liceat, erit \(CB'=BC\mathord{.}\cos\omega\), et proin
\[
\begin{aligned}
AC+CB'&=AC+BC\mathord{.}\cos\omega=AB-BC\mathord{.}(1-\cos\omega)\\
&=AB'-BC(1-\cos\omega)
\end{aligned}
\]
i. e. transitus a puncto \(A\) ad \(B'\) per punctum \(C\) brevior linea brevissima. Q. E. A.
\Article{16}
Theoremati art. praec. associamus aliud, quod ita enunciamus. \textit{Si in superficie curva concipitur linea qualiscunque, a cuius punctis singulis proficiscantur sub angulis rectis et versus eandem plagam innumerae lineae brevissimae aequalis longitudinis, curva, quae earum extremitates alteras iungit, illas singulas sub angulis rectis secabit.} Ad demonstrationem nihil in analysi praecedente mutandum est, nisi quod \(\varphi\) designare debet longitudinem curvae \textit{datae} inde a puncto arbitrario numeratam, aut si mavis functionem huius longitudinis; ita omnia ratiocinia etiamnum valebunt, ea modificatione, quod veritas aequationis \(S=0\) pro \(r=0\) nunc iam in ipsa hypothesi implicatur. Ceterum hoc alterum theorema generalius est praecedente, quod adeo in illo comprehendi censeri potest, dum pro linea data adoptamus circulum infinite parvum circa centrum \(A\) descriptum. Denique monemus, hic quoque considerationes geometricas analyseos vice fungi posse, quibus tamen, quum satis obviae sint, hic non immoramur.
\Article{17}
Revertimur ad formulam \(\BareRoot{(E\mathrm{d}p^2+2F\mathrm{d}p\mathord{.}\mathrm{d}q+G\mathrm{d}q^2)}\), quae indefinite
